\documentclass[../DoAn.tex]{subfiles}
\begin{document}

\begin{center}
    \Large{\textbf{TÓM TẮT ĐỒ ÁN TỐT NGHIỆP}}\\
\end{center}
\vspace{0.5cm}

Ô nhiễm không khí là một trong những thách thức nghiêm trọng nhất về môi trường và sức khỏe cộng đồng tại Việt Nam, tuy nhiên mạng lưới quan trắc PM2.5 mặt đất của quốc gia vẫn còn thưa thớt --- chỉ khoảng 40 trạm quy chuẩn trên diện tích 331.000 km\textsuperscript{2} --- khiến phần lớn dân số thiếu thông tin về chất lượng không khí địa phương. Học máy kết hợp dữ liệu vệ tinh đã được đề xuất như một giải pháp, với các nghiên cứu tại các khu vực có mạng lưới dày đặc thường báo cáo hệ số xác định (R\textsuperscript{2}) trên 0,8. Luận văn này đặt câu hỏi liệu hiệu suất tương đương có thể đạt được cho mạng lưới thưa thớt, nhiệt đới của Việt Nam hay không, đồng thời đối mặt với một vấn đề phương pháp hệ thống: sự khác biệt giữa đánh giá chéo ngẫu nhiên được sử dụng trong hầu hết các nghiên cứu và đánh giá chéo không gian phản ánh nhiệm vụ thực tế của dự đoán tại các vị trí chưa được quan trắc.

Mô hình XGBoost-DART với 66 đặc trưng được huấn luyện trên 37 trạm quy chuẩn, sử dụng độ sâu quang học sol khí từ MODIS và Himawari, các khí vết từ TROPOMI, khí tượng ERA5, sản phẩm mô hình vận chuyển hóa học (CTM) và các yếu tố sử dụng đất. Kết quả trung tâm là chuỗi ba con số: R\textsuperscript{2} xấp xỉ 0,80 với đánh giá chéo ngẫu nhiên giảm xuống xấp xỉ 0,20 với đánh giá LOSO (bỏ-một-trạm-ra), và xuống trung vị khoảng 0,04 cho mỗi trạm khi thông tin oracle bị rút lại. Phân tầng theo mức độ ô nhiễm (T4F) là đòn bẩy hiệu suất lớn nhất, nhưng mang theo sự phụ thuộc vòng tròn --- gán tầng cho trạm đòi hỏi biết trung bình PM2.5, chính là đại lượng cần dự đoán --- và bảy họ phương pháp đại diện đều thất bại trong việc giải quyết nó. Các sản phẩm CTM và vệ tinh toàn cầu thất bại theo các cách riêng biệt: GEOS-CF có chu kỳ ngày đêm đảo ngược và sai lệch trên 200\%, MERRA-2 đạt chỉ số IOA chỉ 0,39, và GHAP chỉ đạt khoảng 50\% độ phù hợp phân tầng. Kiểm chứng ngoại tại cho kết quả tốt trong phạm vi phủ sóng dày đặc, với trung vị R\textsuperscript{2} = 0,53 cho cảm biến chi phí thấp và R\textsuperscript{2} = 0,68 tại Đại sứ quán Hoa Kỳ, giảm dần theo khoảng cách. Ràng buộc chính là mật độ trạm chứ không phải năng lực thuật toán.

\vspace{0.5cm}
\textbf{Từ khóa:} PM2.5, viễn thám, độ sâu quang học sol khí, học máy, đánh giá chéo không gian, Việt Nam, cảm biến chi phí thấp

\vspace{1cm}
Sinh viên thực hiện\\
\vspace{1cm}
[Họ và tên sinh viên]

\end{document}
